\begin{proof}
    First, it must be shown that $( \odot )$ is unital with unit $( -1 )$.
    Since $\Log( -1 ) = i\pi$, then the following equations hold for all $U \in \Mor( \CatUnitary )$.
    \begin{align*}
        (-1) \odot U
        &=
        \exp( \Log(-1) \otimes \Log(U) / (i\pi) )
        =
        \exp( \Log( U ) )
        =
        U
        \\
        U \odot (-1)
        &=
        \exp( \Log(U) \otimes \Log(-1) / (i\pi) )
        =
        \exp( \Log( U ) )
        =
        U
    \end{align*}
    Then $(-1) \odot U = U = U \odot (-1)$ for all $U \in \Mor( \CatUnitary )$.
    In other words, $( \odot )$ is unital with unit $( -1 )$.
    It remains to be shown that $( \odot )$ is associative.
    To this end, let $U, V, W \in \Mor( \CatUnitary )$.
    Let $\{ \alpha_j \}_{j=1}^n$ denote the eigenvalues of $\Log( U )$ with orthonormal eigenbasis $\{ b_j \}_{j=1}^n$ and $\{ \beta_j \}_{j=1}^m$ denote the eigenvalues of $\Log( V )$ with orthonormal eigenbasis $\{ c_j \}_{j=1}^m$.
    Then then following equation holds where $A = \Log( U ) \otimes \Log( V ) / (i\pi)$, since $( \otimes )$ is bilinear.
    \begin{equation*}
        A
        =
        \frac{1}{i\pi}
        \left( \sum_{j=1}^n \alpha_j \ket{b_j} \bra{b_j} \right)
        \otimes
        \left( \sum_{k=1}^m \beta_k \ket{c_k} \bra{c_k} \right)
        =
        \sum_{j=1}^n
        \sum_{k=1}^m
        \frac{\alpha_j\beta_k}{i\pi} ( \ket{b_j} \otimes \ket{c_k} ) ( \bra{b_j} \otimes \bra{c_k} )
    \end{equation*}
    Then $\{ \ket{b_j} \otimes \ket{c_k} : j \in \{ 1, 2, \ldots, n \}, k \in \{ 1, 2, \ldots, m \} \}$ is an orthonormal eigenbasis for $A$.
    Moreover, for each $j \in \{ 1, 2, \ldots, n \}$ and $k \in \{ 1, 2, \ldots, m \}$, the eigenvalue associated with $\ket{b_j} \otimes \ket{c_k}$ is $(\alpha_j\beta_k)/(i\pi)$.
    Since $\beta_k \in i(-\pi,\pi]$, then $\beta_k/(i\pi) \in (-1,1]$.
    Consequently, $(\alpha_j\beta_k)/(i\pi) \in i(-\pi,\pi]$.
    Since $j$ and $k$ were arbitrary, then all eigenvalues of $A$ fall within $i(-\pi,\pi]$.
    Consequently, $\Log( U \odot V ) = A$ and the following equation holds by the associativity and bilinearity of $( \otimes )$.
    \begin{equation*}
        ( U \odot V ) \odot W
        =
        \exp\left( \frac{\Log( U \odot V ) \otimes \Log(W)}{i\pi} \right)
        =
        \exp\left( \frac{ \Log(U) \otimes \Log(V) \otimes \Log(W)}{(i\pi)^2} \right)
    \end{equation*}
    By a similar argument, $\Log( V \odot W ) = B$ where $B = \Log( V ) \otimes \Log( W ) / (i\pi)$.
    Then the following equation also holds by the associativity and bilinearity of $( \otimes )$.
    \begin{equation*}
        U \odot ( V \odot W )
        =
        \exp\left( \frac{\Log(U) \otimes \Log( V \odot W )}{i\pi} \right)
        =
        \exp\left( \frac{ \Log(U) \otimes \Log(V) \otimes \Log(W)}{(i\pi)^2} \right)
    \end{equation*}
    In conclusion, $( U \odot V ) \odot W = U \odot ( V \odot W )$.
    Since $U$, $V$, and $W$ were arbitrary, then $( \odot )$ is associative.
    Next, the equation for $\bigodot_{j=1}^n U_j$ will be derived.
    The proof follows by induction.
    \begin{itemize}
    \item \textbf{Base Case}.
          $\bigodot_{j=1}^0 U_j = -1$ and $\Log( -1 ) = i\pi = i\pi \cdot 1 = i\pi \bigotimes_{j=1}^n \frac{\Log( U_j )}{i\pi}$.
    \item \textbf{Inductive Hypothesis}.
          For some $n \in \mathbb{N}$, if $\{ U_j \}_{j=1}^n$ is a sequence over $\Mor( \CatUnitary )$, then $\Log\left( \bigodot_{j=1}^n U_j \right) = i\pi \bigotimes_{j=1}^n \frac{\Log( U_j )}{i\pi}$.
    \item \textbf{Inductive Step}.
          Assume for some $n \in \mathbb{N}$, if $\{ U_j \}_{j=1}^n$ is a sequence over $\Mor( \CatUnitary )$, then $\Log\left( \bigodot_{j=1}^n U_j \right) = i\pi \bigotimes_{j=1}^n \frac{\Log( U_j )}{i\pi}$.
          Let $\{ U_j \}_{j=1}^{n+1}$ be a sequence over $\Mor( \CatUnitary )$.
          As shown previously, $\Log\left( \bigodot_{j=1}^{n+1} U_j \right) = \Log\left( \bigodot_{j=1}^n U_j \right) \otimes \Log( U_{n+1} ) / (i\pi)$.
          Then by the inductive hypothesis, $\Log\left( \odot_{j=1}^{n+1} U_j \right) = \left( i\pi \bigotimes_{j=1}^n \frac{\Log( U_j )}{i\pi} \right) \otimes \frac{\Log( U_{n+1} )}{i\pi} = i\pi \bigotimes_{j=1}^{n+1} \frac{\Log( U_j )}{i\pi}$.
          Then the inductive step holds.
    \end{itemize}
    Then by the principle of induction, $\Log\left( \bigodot_{j=1}^n \right) = i\pi \bigotimes_{j=1}^n \frac{\Log( U_j )}{i\pi}$ for each sequence $\{ U_j \}_{j=1}^n$ over $\Mor( \CatUnitary )$.
    In other words, $\bigodot_{j=1}^n U_j = \exp\left( i\pi \bigotimes_{j=1}^{n} \frac{\Log( U_j )}{i\pi} \right)$.
\end{proof}